% Corps du contrôle 06 - Proportionnalité (partie 1)

\bareme{
	\exbadge{1}{5} \quad
	\exbadge{2}{4} \quad
	\exbadge{3}{4} \quad
	\exbadge{4}{5} \quad
	\exbadge{5}{2} \quad
	\bonusbadge{2}
}

\begin{consignebox}
	\faIcon{info-circle} \textbf{CONSIGNES IMPORTANTES}
	\begin{itemize}[leftmargin=*,noitemsep,topsep=5pt]
		\item[\faIcon{check}] Répondez directement sur le sujet.
		\item[\faIcon{check}] Justifiez quand c'est demandé (tableau, calcul, phrase).
		\item[\faIcon{check}] Soignez la présentation : unités, phrases, calculs.
	\end{itemize}
\end{consignebox}

\smallspace

\QCMTemplate{QCM — Reconnaître la proportionnalité}{5}{
La situation est-elle proportionnelle ? &
\centering A : Oui &
\centering B : Non &
\centering C : On ne peut pas savoir &
\centering\arraybackslash \solution[1.5cm]{\textbf{A}} \\
\hline
Dans le tableau ci-dessous, les grandeurs sont-elles proportionnelles ?
\[
\begin{array}{|c|c|c|c|}
\hline
2 & 4 & 6 & 8\\
\hline
3 & 6 & 9 & 12\\
\hline
\end{array}
\]
&
\centering A : Oui &
\centering B : Non &
\centering C : On ne peut pas savoir &
\centering\arraybackslash \solution[1.5cm]{\textbf{A}} \\
\hline
Dans le tableau ci-dessous, les grandeurs sont-elles proportionnelles ?
\[
\begin{array}{|c|c|c|c|}
\hline
1 & 2 & 3 & 4\\
\hline
2 & 4 & 7 & 8\\
\hline
\end{array}
\]
&
\centering A : Oui &
\centering B : Non &
\centering C : On ne peut pas savoir &
\centering\arraybackslash \solution[1.5cm]{\textbf{B}} \\
\hline
Un paquet de 5 gommes coûte 2,50 €. Combien coûte 1 gomme ? &
\centering A : 0,20 € &
\centering B : 0,50 € &
\centering C : 1,25 € &
\centering\arraybackslash \solution[1.5cm]{\textbf{B}} \\
\hline
Le coefficient multiplicateur pour passer de 7 à 28 est : &
\centering A : 4 &
\centering B : 21 &
\centering C : 35 &
\centering\arraybackslash \solution[1.5cm]{\textbf{A}} \\
\hline
Dans une situation de proportionnalité, la représentation graphique est : &
\centering A : une droite ne passant pas par l'origine &
\centering B : une droite passant par l'origine &
\centering C : une courbe &
\centering\arraybackslash \solution[1.5cm]{\textbf{B}} \\
\hline
}

\ExTemplate{2}{Tableau de proportionnalité}{4}{table}{
\textbf{1.} Voici un tableau :
\[
\begin{array}{|c|c|c|c|}
\hline
\text{Nombre de cahiers} & 2 & 5 & 8\\
\hline
\text{Prix (en €)} & 3 & 7{,}50 & 12\\
\hline
\end{array}
\]

\begin{enumerate}[label=\textcolor{primary}{\textbf{\alph*)}},leftmargin=*]
	\item Ces grandeurs sont-elles proportionnelles ? Justifier. \points{2}
	
	\anslines{3}{
		Oui : \(\dfrac{3}{2}=\dfrac{7{,}50}{5}=\dfrac{12}{8}=1{,}5\). Le prix est proportionnel au nombre de cahiers.
	}
	
	\item Quel est le coefficient de proportionnalité (prix pour 1 cahier) ? \points{2}
	
	\anslines{2}{
		Coefficient : \(\dfrac{3}{2}=1{,}5\). Un cahier coûte 1{,}50 €.
	}
\end{enumerate}
}

\ExTemplate{3}{Compléter un tableau}{4}{pen}{
On considère le tableau de proportionnalité suivant :
\[
\begin{array}{|c|c|c|c|c|}
\hline
\text{Temps (min)} & 6 & 10 & 15 & 30\\
\hline
\text{Distance (km)} & 1{,}5 & \solution[2cm]{2{,}5} & \solution[2cm]{3{,}75} & \solution[2cm]{7{,}5}\\
\hline
\end{array}
\]

\textbf{1.} Donne le coefficient pour passer des minutes aux kilomètres. \points{2}

\anslines{2}{
	Pour 1 min : \(\dfrac{1{,}5}{6}=0{,}25\) km, donc coefficient \(0{,}25\).
}

\textbf{2.} Explique un calcul permettant de trouver la distance en 30 minutes. \points{2}

\anslines{2}{
	\(30 = 5\times 6\) donc distance \(= 5\times 1{,}5 = 7{,}5\) km (ou \(30\times 0{,}25\)).
}
}

\begin{problembox}{\faIcon{calculator} Exercice 4 : Problème — Passage à l'unité \points{5}}
	Un cycliste parcourt \textbf{3 km en 12 minutes} à vitesse constante.
	
	\begin{enumerate}[label=\textcolor{primary}{\textbf{\alph*)}},leftmargin=*]
		\item Quelle distance parcourt-il en 1 minute ? \points{2}
		
		\anslines{2}{
			En 1 min : \(\dfrac{3}{12} = 0{,}25\) km (soit 250 m).
		}
		
		\item Quelle distance parcourt-il en 45 minutes ? \points{2}
		
		\anslines{2}{
			En 45 min : \(45\times 0{,}25 = 11{,}25\) km.
		}
		
		\item Combien de temps met-il pour parcourir 10 km ? \points{1}
		
		\anslines{1}{
			\(10 \div 0{,}25 = 40\) minutes.
		}
	\end{enumerate}
\end{problembox}

\exspace

\begin{exercisebox}{\faIcon{gift} Exercice 5 : Bonus \points{2}}
	Un paquet de \textbf{4 stylos} coûte \textbf{6 €}. (Situation de proportionnalité.)
	
	\smallspace
	
	Quel est le prix de \textbf{10 stylos} ? \points{2}
	
	\anslines{2}{
		Prix d'un stylo : \(6 \div 4 = 1{,}5\) €. Donc 10 stylos : \(10\times 1{,}5 = 15\) €.
	}
\end{exercisebox}
